\section{Methodological observations}

\begin{frame}{Reproducibility issues uncovered}
\begin{itemize}\setlength\itemsep{4pt}
  \item \textbf{Hard-coded LSTM warm-up.}
        \texttt{init\_period = 252*5} in the downstream LSTM assumes 5 years of trading days.
        Breaks silently for GOOG and ZC$=$F (insufficient history); both run but evaluate the LSTM in its warm-up regime.
  \item \textbf{Device mismatch in \texttt{train\_evolution\_module}.}
        Target tensors are not moved to GPU. Forward stays on CPU even with a CUDA model; training is $8$--$10\times$ slower and the optimiser trajectory differs.
  \item \textbf{$x$-range multiplier in \texttt{plot\_downstream\_tatr}.}
        Augmentation axis multiplied by 10 in plotting; mis-labels ``$10\times$'' as ``$100\times$''.
  \item \textbf{Under-specified synthetic-data protocol} --- the paper does not pin down ``authors'' vs.\ continuous-chain rollout (split).
\end{itemize}
\note{Four issues. The LSTM warm-up bug is silent: nothing crashes, but the model is under-trained on two of three assets. The device-mismatch issue costs an order of magnitude in wall time. The plot multiplier shifts the axis without warning. Together with the protocol question, these four points explain most of the gap between our numbers and the published numbers.}
\end{frame}


\section{Methodological observations}

\begin{frame}{Why exact reproduction was difficult}
\begin{itemize}\setlength\itemsep{5pt}
  \item \textbf{Synthetic-data protocol is under-specified.}
        The paper does not fully specify whether synthetic paths are generated as
        independent segments, one continuous chain, or a continuous chain later split
        into evaluation blocks.

  \item \textbf{Downstream split is not portable across assets.}
        The released downstream LSTM uses \texttt{init\_period = 252*5}.
        This is compatible with the S\&P 500 setup, but GOOG and ZC$=$F do not
        have enough held-out history under the same split. Any valid GOOG/ZC$=$F
        run therefore requires an adapted protocol.

  \item \textbf{GPU training path requires a patch.}
        In \texttt{train\_evolution\_module}, target tensors are not moved to the
        model device. On CUDA this creates a device mismatch, so the intended GPU
        path is not directly reproducible without modifying the code.

  \item \textbf{Published trends are sensitive to protocol choices.}
        Our experiments show that changing the synthetic rollout protocol can
        materially change the TATR trend, even with the same assets and similar
        hyperparameters.
\end{itemize}

\note{
Main message: this is not just a matter of random seeds. The released code and
paper leave several degrees of freedom that materially affect the result. The
most important one is the synthetic generation protocol.
}
\end{frame}


\begin{frame}{Markov-chain stationarity}
\begin{block}{Setup}
PEM defines a Markov chain over $s_m=(p_m,\alpha_m,\beta_m)$. Stationarity would require a fixed invariant $\pi$ with $\pi^\top T = \pi^\top$.
\end{block}
\begin{itemize}\setlength\itemsep{2pt}
  \item Empirical $T$ (per asset, per $K$) is right-stochastic; principal left-eigenvector gives $\pi$.
  \item Stationarity check: simulate long chains, monitor $\|\hat\pi_t-\pi\|_1$ over time.
  \item Drift check: compare mean motif occupancy in real vs.\ synthetic across the chain.
\end{itemize}
\begin{alertblock}{Open question --- empirical burn-in}
Under the \texttt{split} protocol, no clear stationarity is reached within one simulated year. How many segments of burn-in are needed before the chain forgets its init state? Not answered here.
\end{alertblock}
\note{We did produce the MC figures (\texttt{figures/mc\_analysis/}) but did not estimate a burn-in numerically. That is the natural next step.}
\end{frame}

\section{Limitations and Extensions}

\begin{frame}{Limitation 1: motif discovery is fragile}
\framesubtitle{The whole generator depends on the first segmentation step}

\begin{block}{Discrete motif extraction}
SISC turns a continuous time series into a finite dictionary of motifs. This is powerful, but it creates several sensitivities.
\end{block}

\begin{itemize}
\item The number of clusters \(K\) and length range \([l_{\min},l_{\max}]\) are important hyperparameters.
\item Greedy segmentation can misplace boundaries between adjacent motifs.
\item DTW captures shape similarity, but may ignore economically meaningful differences.
\item Motifs are learned offline; new market regimes may require re-clustering.
\end{itemize}

\medskip

\alert{Implication:} generation quality is only as good as the learned motif vocabulary.

\end{frame}


\begin{frame}{Limitation 2: local patterns may miss long-range structure}
\framesubtitle{Short motifs are not the same as full market dynamics}

\begin{columns}[onlytextwidth]
\begin{column}{0.48\textwidth}
\begin{block}{Captured well}
\begin{itemize}
\item Local shape recurrence
\item Heavy-tailed return behavior
\item Volatility clustering proxies
\item Segment-level scaling
\end{itemize}
\end{block}
\end{column}

\begin{column}{0.48\textwidth}
\begin{block}{Harder to capture}
\begin{itemize}
\item Long-memory dependence
\item Regime persistence
\item Macro or news conditioning
\item Rare crises and structural breaks
\item Cross-asset co-movement
\end{itemize}
\end{block}
\end{column}
\end{columns}

\medskip

The learned transition is essentially a compressed state evolution:
\[
(p_m,\alpha_m,\beta_m)
\mapsto
(p_{m+1},\alpha_{m+1},\beta_{m+1}).
\]

\end{frame}


\begin{frame}{Limitation 3: realism does not imply no artificial alpha}
\framesubtitle{Stylized facts are necessary but not sufficient}

The paper evaluates whether synthetic series resemble real data through distributional tests,
stylized facts, and downstream prediction experiments.

\medskip

\begin{block}{Remaining concern}
A synthetic generator can match heavy tails and volatility clustering, while still introducing
\alert{systematic predictable drift} under a simple trading signal.
\end{block}

\begin{itemize}
\item Good marginal fit does not guarantee correct conditional dynamics.
\item Prediction experiments may depend on the model architecture used for evaluation.
\item Leakage is possible if motif extraction or scaling is fitted using future data.
\item A generator useful for augmentation should not create spurious predictability.
\end{itemize}

\end{frame}

\iffalse

\begin{frame}{Extension 1: stricter replication protocol}
\framesubtitle{Make the evaluation time-series safe}

\begin{enumerate}
\item Split the data into chronological windows:
\[
\text{train} \rightarrow \text{validation} \rightarrow \text{final hold-out}.
\]

\item Fit every object on the training window only:
\begin{itemize}
\item motif dictionary;
\item scaling autoencoder;
\item diffusion generator;
\item transition network.
\end{itemize}

\item Generate synthetic paths only after fitting is complete.

\item Evaluate real, synthetic, and null data under the same fixed protocol.
\end{enumerate}

\medskip

\alert{Goal:} avoid using future information when constructing motifs or transition rules.

\end{frame}
\fi

\begin{frame}{Why no predictable drift matters}
\framesubtitle{Martingale intuition for financial returns}

A standard benchmark assumption in financial market modelling is that returns should not contain
an easily exploitable conditional mean under the available information:
\[
\operatorname{E}\!\left[r_{t+1}\mid \mathcal{F}_t\right] = 0,
\]
or, equivalently, for every admissible test function \(g(\mathcal{F}_t)\),
\[
\operatorname{E}\!\left[r_{t+1} g(\mathcal{F}_t)\right] = 0 = \left(
\operatorname{Cov}(r_{t+1},g(\mathcal{F}_t))
+
\operatorname{E}[r_{t+1}]\operatorname{E}[g(\mathcal{F}_t)]
\right)..
\]



\begin{block}{Financial interpretation}
Given the information available at time \(t\), the next return should not be systematically predictable.
Otherwise, the generated data may contain \alert{artificial alpha}.
\end{block}

\begin{itemize}
\item Real financial returns are noisy and close to unpredictable at short horizons.
\item A good synthetic generator should match stylized facts without creating fake trading signals.
\item Therefore, distributional realism alone is not enough: we also need a conditional-mean diagnostic.
\end{itemize}

\end{frame}


\begin{frame}{Extension 1: no-predictable-drift diagnostic}
\framesubtitle{A fixed test for artificial conditional mean}

Use a simple fixed information set:
\[
\omega_t =
\left(1,\ r_t,\ r_{t-1},\ |r_t|,\ |r_{t-1}|\right)^\top .
\]

Define the drift-violation statistic:
\[
\delta
=
\sup_{\|b\|_2 \leq 1}
\left|
\widehat{\operatorname{E}}
\left[
r_{t+1} b^\top \omega_t
\right]
\right|.
\]

For this linear class:
\[
\delta
=
\left\|
\widehat{\operatorname{E}}
\left[
r_{t+1}\omega_t
\right]
\right\|_2 .
\]

\begin{block}{Compare three objects}
\[
\delta_{\text{real}},
\qquad
\delta_{\text{synthetic}},
\qquad
\delta_{\text{null}} .
\]
The null can be built by block-shuffling returns to preserve rough marginal behavior while breaking predictability.
\end{block}

\end{frame}

\begin{frame}{Extension 2: drift-regularized generation}
\framesubtitle{Keep stylized facts, discourage artificial alpha}

Modify training with one additional penalty:
\[
\mathcal{L}_{\text{total}}
=
\mathcal{L}_{\text{FTS}}
+
\lambda
\delta_{\text{synthetic}} .
\]

\begin{itemize}
\item \(\mathcal{L}_{\text{FTS}}\): original reconstruction, diffusion, and transition losses.
\item \(\delta_{\text{synthetic}}\): fixed no-predictable-drift statistic.
\item \(\lambda\): chosen once on validation data.
\end{itemize}

\medskip

\begin{block}{Evaluation question}
Does the penalty reduce predictable drift without destroying:
\begin{itemize}
\item heavy-tailed returns;
\item volatility clustering;
\item distributional similarity to real data?
\end{itemize}
\end{block}

\end{frame}


\begin{frame}{Extension 3: broader model directions}

\begin{columns}[onlytextwidth]
\begin{column}{0.48\textwidth}
\begin{block}{Multivariate FTS-Diffusion}
Generate vector returns:
\[
r_t \in \mathbb{R}^d
\]
and preserve cross-correlations, co-jumps, and sector-level dependence.
\end{block}

\begin{block}{Regime-conditioned transitions}
Condition the transition module on market state:
\[
\phi(p_m,\alpha_m,\beta_m,c_m).
\]
\end{block}
\end{column}

\begin{column}{0.48\textwidth}
\begin{block}{Stress-aware generation}
Oversample rare but realistic regimes: crises, crashes, liquidity shocks, and volatility spikes.
\end{block}

\begin{block}{Stronger validation}
Move beyond distribution matching:
\begin{itemize}
\item martingale-style diagnostics;
\item trading-rule robustness;
\item risk metrics such as VaR and ES;
\item out-of-sample transfer to new assets.
\end{itemize}
\end{block}
\end{column}
\end{columns}

\end{frame}


\begin{frame}{Takeaway}

\begin{block}{Main limitation}
FTS-Diffusion is strong at generating realistic local motif-based financial paths, but its guarantees are mostly empirical and distributional.
\end{block}

\begin{block}{Main extension}
Add a fixed no-predictable-drift diagnostic and, optionally, a drift penalty:
\[
\text{realistic synthetic data}
\quad + \quad
\text{no artificial predictability}.
\]
\end{block}

\medskip

This gives a cleaner benchmark for deciding whether synthetic financial time series are useful for augmentation rather than merely visually or marginally realistic.

\end{frame}



% --- Slide: Conclusion ---
\begin{frame}{Conclusion}
    \begin{itemize}
        \item FTS-Diffusion addresses the core challenges of irregularity and scale-invariance in financial data.
        \item SISC allows for clustering recurring market ``shapes'' effectively.
        \item The model demonstrates superior performance in preserving statistical properties compared to TimeGAN and QuantGAN.
    \end{itemize}
\end{frame}


